\section{Western Social Order}

\paragraph{Part I} Cologero points out that quibbling over minor points of philosophy \& actualizing states of being are not equivalent for the noble character. That is, philosophical debate is not for the gentleman, beyond a certain point. Surely the West (as such) was built upon such a a fundamental impulse, as the aristocracy of the Franks and Germans were a rough and ready group of fighters, with little taste for Byzantine theology. John Romanides\footnote{\url{http://www.romanity.org/htm/rom.03.en.franks_romans_feudalism_and_doctrine.01.htm}} has criticized this substantially, but does not address in any way the fact that the West managed to not merely retain much more than a shadow of Byzantine mysticism, but actually to incorporate and seemingly ``add'' elements of Tradition with which the East was less familiar. Cologero has discussed this under the title of The Three Orders\footnote{\url{http://www.gornahoor.net/?p=5198}}. Byzantium retained the purity of dogma from Constantine's day, but the social order was collapsing. They were apparently unable to escape the curse of political factions: in other words, they lacked unity\footnote{\url{http://en.wikipedia.org/wiki/Byzantine_Greeks}}. Procopius\footnote{\url{http://en.wikipedia.org/wiki/Procopius}} even relates a story about Justinian the emperor which is intriguing;

\begin{quotex}
And some of those who have been with Justinian at the palace late at night, men who were pure of spirit, have thought they saw a strange demoniac form taking his place. One man said that the Emperor suddenly rose from his throne and walked about, and indeed he was never wont to remain sitting for long, and immediately Justinian's head vanished, while the rest of his body seemed to ebb and flow; whereat the beholder stood aghast and fearful, wondering if his eyes were deceiving him. But presently he perceived the vanished head filling out and joining the body again as strangely as it had left it
\end{quotex}


... perhaps some of the ``pure in spirit'' who could perceive his disappearing head were Pythagoreans. In any case, political factions dismembered Byzantium long before the Seljuk Turks and the Crusaders delivered the finishing blows. Although I don't agree with the pejorative, the adjective ``Byzantium'' today still describes a certain kind of stifling, convoluted atmosphere that is almost impossible to fathom. In fact, Cologero has also pointed out that the East/West schism was primarily exoteric, \& should not (in fact) be recognized as decisive or definitive by those who have eyes to see. Just as the Inquisition's attack on esoteric bodies within the Church should be ``ignored'', in the sense of recognizing it as a natural kind of failing in these situations, while retaining both facets for future use, the West vs. East problem is also illusory, if one is trying to deduce eternal principles out of historical dialectic. Instead, move from Unity towards the One, rather than reasoning from the many back to Unity.

The Revolution is derived from such dialectic; although Romanides is not a revolutionary, his theology has made it in some senses more difficult to rapproche with West, in that there is nothing constructive or creative about it, and one has to ``supply'' what is missing. Indeed, had it not been for Gornahoor, I myself would have found his logic convincing, and would no doubt by now be reading Alexander Dugin and plotting to immigrate to Russia, or at least pining for it. A Revolutionary reasons thusly:

1. The Church has done bad things, or been implicated in the doing of said bad things.

2. Therefore, we are better off today with the Church in chains, culturally, if not annihilated forever.

Therefore, the upshot today is that, ``as for my people, women are their oppressors, and children rule over them''. Isaiah 3:12.

As Robert Nisbet has pointed out in \emph{Twilight of Authority}\footnote{\url{http://www.amazon.com/Twilight-Authority-Robert-Nisbet/dp/0865972125}}, the modern liberal state that has resulted from the above reasoning (which is shared by some of the brightest and best young people I have met or been acquainted with) drifts toward either chaos or a monolith that reduces individuals to atoms, with no power or defense against the centralized state, which has been optimized (ostensibly) for the good of all, and particularly the new ``individual'' neo-bourgeois\footnote{\url{http://www.firstthings.com/web-exclusives/2012/08/the-neo-bourgeois-project}}. Phillip Rieff is a Jew\footnote{\url{http://en.wikipedia.org/wiki/Philip_Rieff}} that every man of Tradition ought to read: his book on \emph{Teaching \& The Second Death} closed off, forever for me, the idea that the new ``individual'' had anything to add -- his diagnosis of the spiritual sickness that grips us is profound. It is no use arguing, you have to treat the condition as an illness; in this Romanides and Tomberg are in essential agreement, as are all the Traditional thinkers.

Unfortunately, a lot of the young are ineluctably drawn to the syllogism above, and those of old age, who are hopelessly lazy and self-corrupt, encourage them. What the enthymeme\footnote{\url{http://en.wikipedia.org/wiki/Enthymeme}} above leaves out is significant. What, for example, happens to other Ideals besides the Church that become tarnished? Are they to be discarded as well? And what happens when Man has tarnished his last Ideal? Should Ideals themselves be given up? What would that look like? Can Ideals be rehabilitated? Why should the abuse of an Ideal disavow the goodness or reality of that Ideal? Are Ideals inevitable, even when denied? Why did the Ideals fail to begin with? Are they inherently evil, inviting abuse? Why should this be? Can Ideals be critiqued by anything else other than Ideals? Were the Ideals ever properly instituted in the first place? Why or why not? What would a false Ideal look like, and how would we recognize it?

Cologero rightly points out that Socratic dialogue can end in \emph{aporia}. But what if the Socratic dialogue was done, as he suggests, within a framework of actualization, as acts achieving states of being? That is the project of Gornahoor, or at least, the part I am most familiar with and best understand. It may be preaching to the choir, but it's certainly better than what you'll get at the average Sunday School, and you'll know something more than a four-part Gospel harmony, although such things have value.

As an aside of interest, an author that links to the site frequently, Brett Stevens, has written a popular piece\footnote{\url{http://www.amerika.org/politics/destroy-white-nationalism/}} which summarizes in popular form some of the cogent criticisms of political European nationalism that Cologero has made in much more in depth and subtle terms. Some of the readers may find the piece useful and convincing, as I did. It points us in the direction of achieving what our director and founder a scale of spiritual order, rather than appealing to the baser parts of human nature, our own and others.

The West, historically, has excelled at achieving difficult balances with periods of crisis, from its inception, which involved multitudes of groups and periods of disorder. It can be hoped that it will do so again.

\paragraph{Part II} A recent comment by John of Salisbury made me think:

\begin{quotex}
Not only were these scholars unable to drive out the bad scholars, but in combating insanity, they temporarily became insane...
\end{quotex}

I thought about this question in relation to other questions or comments relating to the ``inner struggle'' of the man of Tradition in today's very untraditional world. I had a good acquaintance on the Internet give up study in that area because he needed something more cheerful. If you will notice, Cologero does not restrict his reading to Guenon and Evola, but also studies a wider range of writers on a diverse array of subjects. It is probably true that a steady diet of Evola and Guenon has the potential to unsettle one. Caleb is planning to put up a post that addresses part of this, in that Tomberg's \emph{Meditations} are designed to lay a firm foundation, without which (if you neglect the first card of the Tarot), the rest becomes unstable. Cologero has done a fine job of supplementing Evola and Guenon (whom he terms the ``Master of Tradition''), with the meditations on the \emph{Meditation}. The Tarot discussion list (and Caleb's article) is designed to assist in this process. In trying to imitate our betters, we often overstep or over-estimate our own capabilities, sometimes not deliberately, as we try to ``come up to speed''. Incidentally, much of Phillip Rieff's thought in \emph{Culture and Its Second Death} explores the danger of being initiated into our elder's/better's debates too swiftly, including the fundamentally erotic desire to criticize those to whom we owe everything. Rosenstock-Huessy called this ``Teaching too Early, Learning to Late'', in which the danger is that we learn just enough to be dangerous, to ourselves and others.

\begin{quotex}
A little knowledge is a dangerous thing -- drink deep, or taste not the Pierian spring!
\end{quotex}


I can certainly relate to this, having passed in and through and back into such phases, over the course of my short 40 years. In the case of Gornahoor, the best possible foundation has been laid down, but it has been made freely and openly available (with no secrets) laying out the Arcanum that must be mastered before entrance into the Mysteries. But the danger is still there, just a necessary one: in our chaotic age, one has to invite everyone, since those who should be ready are not. Jesus gave this in parable form in the tale of the Wedding Feast invitations.

John of Salisbury cites a medieval authority who claimed that there are three types of souls: those who fly, those who crawl, those who walk. Those who fly learn much, and swiftly, but soon forget. They make poor students (I think this is my soul's danger, for instance). Poetry may be a help to such students. Those who crawl are dull of soul -- in Pythagorean circles, they would not be allowed to be initiated, for instance, during the old days. Those who walk, steadily and constantly, those who are pilgrims, make the best students, fit to climb Parnassus and drink from its springs.

There was another recent article on Gornahoor that explained some of this quite well -- different types or souls have to be handled\footnote{\url{http://www.gornahoor.net/?p=6980}} differently. Michel's comment on this post\footnote{\url{http://www.gornahoor.net/?p=7093}} by Cologero gets to an esoteric root that lies behind a lot of misunderstandings that occur during these confusing times. Different types or \emph{differing potentials} require completely different spiritual trajectories: in Mouravieff's words, if the center is God, and we all find ourselves scattered at various circumferences, the trajectory back of necessity takes a variety of angles and forms. And some people don't necessarily manage to make a straight line, either!

All of this is to say that the ``care of souls'' remains paramount: even if one is caring for one's own soul, in which case, even more ``care'' is required, not less.

\begin{quotex}
In some denominations of Christianity\footnote{\url{http://en.wikipedia.org/wiki/Christianity}}, the cure of souls (Latin\footnote{\url{http://en.wikipedia.org/wiki/Latin_language}}: \textbf{\emph{cura animarum}}), an archaic translation which is better rendered today as ``care of souls'' is the exercise by priests\footnote{\url{http://en.wikipedia.org/wiki/Priest}} of their office. This typically embraces instruction, by sermons\footnote{\url{http://en.wikipedia.org/wiki/Sermon}}, admonitions and administration of sacraments\footnote{\url{http://en.wikipedia.org/wiki/Sacrament}}, to the congregation\footnote{\url{http://en.wiktionary.org/wiki/congregation}} over which they have authority from the church\footnote{\url{http://en.wikipedia.org/wiki/Christian_Church}}. In countries where the Roman Catholic Church\footnote{\url{http://en.wikipedia.org/wiki/Roman_Catholic_Church}} acted as the national church, the ``cure'' was not only over a congregation or congregations, but over a district. The assignment of a priest to a district subdividing a diocese was a process begun in the 4th century AD. The term \emph{parish}\footnote{\url{http://en.wikipedia.org/wiki/Parish}} as applied to this district comes from the Greek word for district. Those who earned their living on a position without cure of souls were known to have a sinecure\footnote{\url{http://en.wikipedia.org/wiki/Sinecure}} (hence the expression).
\end{quotex}


Even when one is taking the ``dry path'', and investigating the layers of one's own soul in an esoteric manner, there is a care or ``art'' involved, \& note merely a science (this may be where the balance between East and West is reached, as the East approaches spiritual things almost mechanistically).

This pervasive feeling of hopelessness, of fighting against fearful odds, of being doomed, can be viewed and attacked from several angles. Since I've been asked several questions lately that pertain to it, I assume it's relevant, and am relating it back to the conquest or reconquest of Western Social Order. It may relate to the ``noonday demon'' that attacked the desert monks, a lassitude and despair that could only be overcome with productive work: this is why Cologero has emphasized that practicing one's caste and craft is as important as becoming ``enlightened''. This is in the spirit of the Patristic fathers, who had monks weave baskets, plant gardens, build buildings, etc., because (as the \emph{Philokalia} teaches) Satan can only attack along one front when we are working, rather than all fronts simultaneously.

First, as noted above, it is paramount to appreciate that simultaneously, Guenon is both the measuring rod (for several reasons, not least of which is that our teacher, Cologero, regards him as such) \& also that which is to be transcended. That is, the goal is to not make Guenon's mistakes. For a Westerner, then, we recognize that Guenon advocates a path away from exoteric form, in favor of an East that (in his day) was less tainted, and a source of hope. Of course, today, we see that the East has imitated the West, and that Coca-Cola and blue jeans have flattened the globe. Cologero has done a huge service in upholding Guenon's standards, but eschewing some of his conclusions.

Secondly, we avoid the pit of getting caught up in erotic disputations over nothings. We've seen a lot of posters come and go -- many wanted to dispute with Cologero over various bagatelles or supposed critical mistakes. The Hermetic method (favored by Western thought that is faithful to Plato, Bonaventura, etc.) eschews such disputes, as does St. Paul. Avoid fruitless disputation... .

Thirdly, for the above reasons, we should probably avoid a pure diet of Guenon and Evola and traditional thinkers, if for no other reason (this is an additional one) than that those who battle vampires and monsters develop a kind of hardness about them that can be monstrous. It is not natural (and they would have recognized this) to have to spend your whole life in combating insanity, from dusk till dawn. Rejoice over their gift to us, but don't forget to plant a fig tree. Our times are conditioned by different possibilities; many battles are lost, others are looming, and some new possibilities have deepened or awakened. As our master Goethe says, \emph{a man who is not of his own times, can be a man of no other time}. Read a mystery novel, watch a movie. Even the medieval thinkers (such as the author of \emph{Philobiblikon}) says that it is possible to appreciate modern thinkers and writers, along with the ancient ones. Evola and Guenon have modern dimensions to their thinking. \textbf{\emph{The entire Quantity/Quality distinction can become (at times) almost a quantitative category which taints our perception of Quality}}. The medievals emphasized the value of Estimation, in which the beginner practiced his craft of working by esteeming that which had objective, natural moral value, in their own day. Even in our day, there are people which have dignity, things or activities which have value, etc.

So temper it. If it gets you down so much you can't function at all, and quenches the inner fire, even down to the last spark, find something more cheerful for awhile. But realize it isn't Guenon or Evola's fault -- its a privation within ourselves. Luckily, within the Christian esoteric and exoteric tradition, we are allowed holy-days and feast days, in which we can re-create ourselves in the wonder of Creation, which abides, even in our horribly dark days. As our teacher Wordsworth says, \emph{one impulse from a vernal wood, can teach you more of man, of moral evil and of good, than all the sages can}. Presumably, Wordsworth didn't even have access to the sages which we do, so temper this with grateful knowledge, and spend some time learning the constellations, or admiring the wonder of a storm. In the \emph{Romance of the Rose}, an initiatic text, the writer asks rhetorically, what can be done with someone who doesn't appreciate the approach of Spring and the gods of Love!

Fourthly, there are other thinkers out there, perhaps secular ones, who can be as circumstantially valuable to us as Evola and Guenon, under certain conditions. I would class Phillip Rieff in this category, at least for those like myself. He gives a feeling of ``utter hopelessness'', but this is meant to be curative, in the spirit of Kafka\footnote{\url{http://books.google.com/books?id=45d6JsdVFaIC\&pg=PA50\&lpg=PA50\&dq=kafka+withered+hand\&source=bl\&ots=yOFVlpXRLO\&sig=VxXgUnJ-qHSfTDsw6og2NQcl8w8\&hl=en\&sa=X\&ei=PL0LU-bSK8XJsQT5vIG4BQ\&ved=0CGkQ6AEwCQ\#v=onepage\&q=kafka\%20withered\%20hand\&f=false}}. And it is aimed at the \emph{hubris} and ignorance of our modern self. Some self-hatred and the cutting off of limbs or eyes is appropriate (symbolically speaking) as most moderns are in a rather unique bind, being cut off from their own soul, which is
\emph{naturaliter anima Christiana} (Tertullian). So a little Theodore

Dalrymple, Richard Weaver, Phillip Rieff, Christopher Lasch, or other such notable secular saints is of great value to many -- if nothing else, it's a safer way to process (like the liver) the toxins of our intellectual age, which are in the air we breathe. Just realize that (being creatures) we have to balance our diet intellectually (until we are enlightened, when all things become pure). \emph{If you read Rieff, you'll realize you should probably be listening to Haydn's classical music (or some other classical composers!)}. However, the trip is worth it. Rieff shows us (in \emph{Culture, Its Second Death}) that all of our spiritual states (even the Sartrean hell, which lurks for us today) are manifestations of our soul's possibilities along the Vertical. There is no escape (or no false escape) from the vertical -- we have to face up to the real problems and questions, first. As Solomon says, the house of mourning is more fruitful than the house of laughter, if it is a mourning of repentance. If (however) one is past this, then one is past it.

Christians are ``ahead of the times'': we are not immersed in them. So just because this is the Kali Yuga, it (everything) still matters very much. One of my teachers, a Lithuanian who taught literature, once berated me for reading Christianity into Shalamov's\footnote{\url{http://en.wikipedia.org/wiki/Varlam_Shalamov}} work: the whole point of his story was how small acts had meaning in and of themselves, regardless of a Christian interpretation or outcome! The character (interned in a camp) had found a can of condensed milk. For a short time, this can was God's grace, without being officially God's grace! Such is the mighty, eternal, and infinite mercy of God, which mercy endures forever, His chief attribute. This is the lesson Tomberg teaches also -- not control, but mercy and Love, which is the Queen of Magic.

This is our burden, this is our time. It is hard, and we have to help each other. We have to show mercy, while remembering Truth. I hope this is encouraging, \& look forward to being taught and helped by others on this path.

\paragraph{Part III} In the \emph{Philobiblon\footnote{\url{http://www.philobiblon.com/philobiblon.shtml}}}, Richard de Bury (Bishop of Durham) justifies his love of books in a Christian framework by pointing out that Plato is said to have paid 10000 dinars for a rare scroll of Philolaus. Philolaus'\footnote{\url{http://en.wikipedia.org/wiki/Philolaus}} work was most likely a transmission of the teachings of Pythagoras, who had gone to Tyre and to Egypt in order to learn from the masters there. Plato used this work to compose the \emph{Timaeus}, one of the more ``mystical'' of his treatises. If you are wondering where the Eastern Church got its teaching on the total dissimilarity between the essence of God (unknowable) and His energies (especially those manifested in Creation), this is where it comes from: it was part of the Tradition embodied in Hellenic wisdom, culled from Egypt:

\begin{quotex}
This is the state of affairs about nature and harmony. The essence of things is eternal; it is a unique and divine nature, the knowledge of which does not belong to man. Still it would not be possible that any of the things that are, and are known by us, should arrive to our knowledge, if this essence was not the internal foundation of the principles of which the world was founded, that is, of the limiting and unlimited elements. Now since these principles are not mutually similar, neither of similar nature, it would be impossible that the order of the world should have been formed by them, \textbf{unless the harmony intervened} . . .
\flright{\textsc{Philolaus}, \emph{Fragment DK 44B 6a.}}
\end{quotex}

\emph{Out of Egypt}, said the Patristic Fathers, \emph{have I called my Son}. Here was an example of the higher, ``secret'' meaning of Scripture, which they discerned in the literal text, and defended in their apologetics. The \textbf{\emph{Logos}} (of course) is the intervening harmony which causes the co-inherence\footnote{\url{http://web.sbu.edu/friedsam/inklings/coinheretance.htm}} of the invisible higher world (which is Divine, and completely dissimilar) and the mortal, material world (which nevertheless is made ``in the Image''). Through the Logos, argued Saint Paul, the worlds were brought into Being, and they will be recapitulated or ``summed up in Him'' again at the end of Time (Colossians). The Christian ought to be ``ahead of the Times'', because the Logos has in these latter days revealed Himself as co-substantial with the Father, our older brother, a person Who stooped to conquer Death.

The modern pagan rejection of Christianity puts the West in an interesting position. It is a sort of Mexican stand off. On the one hand, the pagans of the classical world actually were incorporated into the Medieval Tradition (arguably with historic difficulties or distortions), and converted to the Church in droves, with notable exceptions. So the stream of spirit moved into Christian civilization, presenting two problems: 1) Some pagan thought was lost or misunderstood, but to the extent it was not, it now presents as ``Christian'' (eg., the doctrine of Purgatory). 2) Tradition now possesses a distinctly Christian cast in the West. It is for these reasons (among others) that Gornahoor has encouraged a revival of Tradition that is a return to the Medieval outlook, both because a rejection of the Logos understood as Christ would be to repeat on a grand scale the historic distortions which did occur when Christianity assimilated Greece \& Rome, and because to an enormous and little understood degree, Christianity was actually meant to be the perfecting and crown of those very pagan mysteries. There are other cogent reasons, more practical than that, as well. Cologero's quote from Tomberg\footnote{\url{http://www.gornahoor.net/?p=7147}} about ``salvation'' is relevant here:

\begin{quotex}
Salvation is the helping hand offered to idealistically minded human beings who, though no longer identifying themselves with natural evolution, still lack sufficient strength of themselves to alter its course.
\end{quotex}


It is perilous, possibly foolhardy, to reject a prophet, vehicle, or messenger, because why would God send another one? Or, put even more precisely (since God's mercy is infinite), why would we recognize the Divine the second time, if we rejected it the first time around?

These, however, are merely practical considerations of various orders, which can always be contested empirically or historically in the form of interminable debate. What is more importantly at stake is the very concept of Logos itself, both metaphysically and practically. Every Tradition had a Tao or Logos which represented the work of the best minds and spirits in apprehending what the Spirit was trying to teach men over vast bodies of space and Time, and the infinite abyss of the Fall (even worse). A rejection of Logos might well be accompanied by a turning inward, like the crawfish in Tomberg's Tarot card, which represents the darkening of the Nous in self-imposed, invincible ignorance.

Since obeying the Logos, then enlightening the Nous (Law and then Grace) represent only initial steps, the first stages of divinization, what is being lost is the very possibility of even preparing for Enlightenment, and even more opportunities wasted in what is a long, long journey.
\textbf{This, in very fact, is exactly what political, social, and spiritual Revolution represent}. Every traditional writer popular with even the New Age crowd (Mouravieff, Tomberg, and Heindel) mention in passing asides the obvious fact that Revolution represents emotional immaturity and darkened Reason, and is an insufferable barrier to progress in the esoteric art. The idea is not that the esoteric man will be ``free'' in the sense that he will never darken a Church door again in his life, but rather that the more ``free'' one becomes, the more one serves and is able to properly fulfill even exoteric duties. \emph{Mouravieff, I believe, mentions that one must train one's replacement, before the Universe will allow one to ``escape'' one's current status quo}. Predictably, this was faithfully reflected in Cologero's meditations, and that portion of Gornahoor's program was just as predictably disliked by certain quarters.

Tomberg lays particularly heavy emphasis both on Love as being \emph{more} perfect than the Eastern ideal of returning to the purely embryonic primordial state (he does not however, deny their validity), and the practice of mercy and service as being integrally bound up with the ``gift of tears'' given to the West. He teaches that the lost sheep of the ``personality'' was destined for salvation, and pleads with the Lord of mercy to forgive him for retaining the use of the human reason and its chakra in his \emph{Meditations: the French Hermetic Christian tradition of Magic was uniquely both human and Christian in this way.} Like Abraham pleading with God for the life of the city, Tomberg suspected that this was God's divine intention all along, and that earlier Traditions were only ``imperfect'' in that they were capable of being preserved, refined, and transformed in a non-Revolutionary way into something even more themselves. That was why he dared to wrestle with the angels. In chapter 2, he even asserts that Love is prior to Being.

And in fact this was how the Christian claim should have been made, and occasionally, was made in the proper manner by very august characters and intellects. There was actually a small revival of it at Cambridge\footnote{\url{http://plato.stanford.edu/entries/cambridge-platonists/}} in old England, but these are far from the best expositors. A better example of it would be Bonaventura, or John Scotus Erigena. Yet the ``baptism'' of what was true, beautiful, and good was certainly how the medieval instinct understood itself to be working, although obviously it was imperfect as it played out in time and space.

Or is it? Unused to seeing Tradition embodied at all, \& undiscerning (in general) of possibilities and limits, the modern world wears its spectacles from a supposed Olympian vantage point in order to condemn the old orders. A Time \& Space in particular (400-1500) isn't a rarefied time and space, but falls in between 700BC-399AD \& 1501-2014 AD. Which means that criticism aimed at this period would reflect upon the periods immediately proceeding and following; to truly judge the period, one would have to resort to Transcendence, and this judgement is much more favorable.

Therefore, Cologero has asserted the distinction between Revolution and ``the opposite of Revolution'', or between Becoming and an Order based on Love's generation of Being. So that we ought to take a different and more charitable view\footnote{\url{http://www.millinerd.com/2013/11/modernity-middle-ages.html}} of the Middle Ages:

\begin{quotex}
Paradox... is not dialectical. Paradox is the simultaneous assertion (not the reconciliation) of opposites. Because of the paradox not just of Christ's incarnation (God in the human) but also of divine creation (God's presence in all that is infinitely distant from him), \textbf{matter was that which both threatened and offered salvation}. It threatened salvation because it was that which changed. But it was also the place of salvation, and it manifested this exactly through the capacity for change implanted in it. When wood or wafer bled, matter showed itself as transcending, exactly by expressing, its own materiality. It manifested enduring life (continuity, existence) in death (discontinuity, rupture, change). Miraculous matter was simultaneously -- hence paradoxically -- the changeable stuff of not-God and the locus of a God revealed (34-35) In this sense, medieval history can never be properly ``understood'' as much as experienced\footnote{\url{http://ayjay.tumblr.com/post/67747782276/if-the-reader-will-suspend-his-disbelief-and}}.
\end{quotex}


Bynum continues:

\begin{quotex}
Paradox is by definition impossible to explain in discursive language. One cannot simultaneously assert contraries. Rude, other-denying facts such as identity and annihilation, or the haunting presence and yet utter beyondness of ultimate meaning, cannot be spoken together. Yet together they must be lived. Their simultaneity cannot be stated; it can only be evoked -- and even this only inadequately.
\end{quotex}


It can easily be seen that a man of Tradition would have an easier time soaking in thought that is affinitive to him by taking the same non-Revolutionary approach common to the classical \& medieval worlds. The Church gets blamed for burning books of pagan knowledge, but it needs to be remembered that a great deal of ancient learning was in fact saved by the Semitic ``axial religions'', baptized by the saints (St. Patrick permitted ``finger magic'' based on the oghams), \& passed down in the common culture of a thoroughly medieval and Christian people. And this does not even begin to address the esoteric or scholarly side of it, which was much more aggressive and open to inheriting the pagan tradition, including embodying it in the Eucharist\footnote{\url{http://deepblue.lib.umich.edu/bitstream/handle/2027.42/64788/jasonbp_1.pdf}}.

We distinguish, therefore, between the occasional or even sustained rhetoric of the Christian authors (including that arch-pagan and super-Christian, Tertullian) and the actual action of the Spirit, reflected in their thought, which actions blows where it wills \& ignites high intellects, pure characters, \& rare hearts in a variety of Traditions. What is the purpose of this variety except to reach those \emph{within} a specific variety, \& to convince those who are rising above it with the use of it, with even greater proofs, that their path is blessed? For since our access is not direct, we approach this through a desire of imitating that initial inspiration, but the fact of inheriting very specific traditions from those above us in the Great Chain of Being. And since it stands to strong Reason that if one is ``risen'' above one's Western tradition, it would seem to serve but little to the purpose to immerse one's self again in the elementary and rudimentary elements, we conclude that the medieval heritage remains of great, even inestimable, value, most particularly for those who desire to rise above it.

Or, as Saint Paul put it,

\begin{quotex}
2 But is under \textbf{tutors} and \textbf{governors} until the time appointed of the father.  3 Even so we, when we were children, were in bondage under the elements of the world:  4 But when the fulness of the time was come , God sent forth his Son, made of a woman, made under the law, 5 To redeem them that were under the law, that we might receive the adoption of sons. 6 And because ye are sons, God hath sent forth the Spirit of his Son into your hearts, crying , Abba, Father. 7 Wherefore thou art no more a servant, but a son; and if a son, then an heir of God through Christ.  8 Howbeit then , when ye knew not God, ye did service unto them which by nature are no gods. 9 But now, after that ye have known God, or rather are known of God, how turn ye again to the weak and beggarly elements, whereunto ye desire again to be in bondage ?  10 Ye observe days, and months, and times, and years. 11 I am afraid of you, lest I have bestowed upon you labour in vain. 
\end{quotex}


If Christianity, too, has a tutelary function in the exoteric realm, it still makes seemingly little sense to surpass one tutor, and to graduate to another. After a while, might one not wonder if the problem was in the pupil, rather than the tutor? Be sure that one is truly seeking the Father, rather than merely desiring to escape the General Law of the Tutor, \& this requires great, even supernatural honesty with the Self.

The emphasis on decency and good order, even in spiritual matters, is characteristic of the medieval outlook. We could do well to profit from it. It accords greatly with being men who are affiliated with that which is the ``opposite of a Revolution'': men of Logos, men of the West.


\flrightit{Posted on 2014-02-14 by Logres}

\begin{center}* * *\end{center}

\begin{footnotesize}
\begin{sffamily}
\texttt{Ash on 2014-02-14 at 19:18 said:}

I am currently reading William Thomas Walsh's work in the Inquisition of Europe. He argues that the dominant heresies of the time stemmed from the Manichees, among others, and represented political as well as spiritual threats to the Christian order, stoked on by conversos and those more loyal to Kings than Pope. The lesson is, there are few who truly view their actions as being ``bad'', neither inquisitors nor heretics nor schismatics. So extending a true higher unity to the social plane is the most arduous of tasks, even in Christendom. Perhaps the answer to the mindset you described is the Christian virtue of charity, even in the harshest of settings like that of a corrupt, failing, or fanatic church.

The reading circle I am currently part of has in its numbers Ismaelis, Westerners outside any Churches like myself, and some Christians. It works when a small number respects that higher unity. We might move from the unity to the One someday. The elite from different faiths may be able to have that unity as well. It will be a miracle indeed if we see it extended to whatever new order arises.

\hfill

\texttt{Logres on 2014-02-14 at 23:46 said:}

Charity would be a much better watchword than tolerance, that is certain. Some good thoughts there, yes it is arduous; is it any wonder there were problems? So many people think the answer is to just give up on the Ideals; I can't see this, as it is a logical and metaphysical impossibility anyway. If our social order reflects our personal orders, then we truly are faced with a big task. Not impossible, however. Not impossible.

\hfill

\texttt{Logres on 2014-02-24 at 22:34 said:}

I was thinking about Phillip Rieff's vertical/horizontal distinction. It occured to me that the Christian cross has a raised horizontal bar, to remind us that the vertical lifts up the horizontal, but that the vertical bar is longest, so that it retains its pre-eminence, even as it serves the material plane. So the Christian cross is modified from the ancient one in a subtle way. Rieff argues that the safest course of life is to live in movement more or less horizontal, but reminds us that this is ``that which ought not to be done, but is permitted'' : this may have an analogy with exoteric religion, which has certain things permitted, but not encouraged. Rieff also insists that there is a difference between transgressive (Revolutionary) order or disorder, and a permissive order (which openly specializes in the horizontal), and a further distinction still between older orders, which merely permitted without dignifying the permission. So his sociology may have some relevance to students of traditional order, like the Middle Ages... 

\hfill

\texttt{Max on 2014-02-25 at 15:08 said:}

The sun sets in the west and while dark continues its transition to north before being born again. Transformation takes place in the dark says Guenon, the highest is reflected in the lowest. How can one navigate the dark landscape without having knowledge of good and evil, what is up or down. We recognize the struggle in ourselves. To flee to some supposed safe haven is no answer, it means death. We accept the challenge to walk through hell. There is no escape anymore. Western man is not one to turn away. The warrior stares the danger straight in its eye and demands; give me all you got! I am still here. I am not afraid.

\hfill

\texttt{JA on 2014-02-26 at 08:28 said:}

I think that the problem with reading solely Guenon and Evola is that they define what is wrong more than propose what should be... ... ... ... ... .

\hfill

\texttt{argusandphoenix on 2014-03-17 at 23:43 said:}

Rieff made the distinction, I was inferring from that insight. A revolutionary order is one that transgresses by denying that there is such a thing as Transgression, but then schizophrenically enforces this through making it a transgression to punish transgressions. Hence, the hatred for the natural moral aristocrat, or otherwise. A permissive order specializes in evading vertical claims by maneuvering with infinite cunning back and forth along the horizontal incline. It negotiates with what is eternal. This is the ``natural state'' of most men. Older orders (and here I might disagree, but think perhaps I am just not understanding) are normative in that they do not attempt to abolish General Law, but rather, sublimate it. They would make room for permission (eg., the brothel and red light district by the railroad tracks) but would officially condemn it. Their attitude towards the transgressors was more intolerant, and can only be understood fully in light of later events (the Red Terror, and such) which came out of the triumph of the inverted order. They insisted on accepting the transgressors if and only if some use or place could be found from them in the existing order (eg., the impressing of criminals into the army, etc.). Otherwise, they left them alone with periodic repressions. I agree that feudalism can't return, per se, but am not sure we don't have it already. What else but a serf is a Walmart employee? Rieff thought the beginning of wisdom was in ``Thou Shalt Not'' and the fear of the Lord. This is the triumph of the vertical that simply can't be surpassed (as I see it, although it can be embodied to some extent, perhaps to a great extent) in the evolution of Time and Space in a literal fashion. Your vision is keen indeed, but few attain the dispassionate effort to see it. What it will look like, from the outside, is rather a slowly ascending spiral, with many more Dark and Golden Ages to come, and more versions of feudalism, revolution, restoration, order. From the inside is a different matter. But history is usually written from the outside. -- Logres

\hfill

\texttt{Pickman on 2014-04-04 at 18:22 said:}

The Imperium is already here, Tradition as it is the cosmos and universal will of the logos is always upon us and acting even on modernity. Far future or past, it's with us now, in the present we simply lack the multidimensional vibrational consciousness to witness it. What is time? What is space? What is materiality -- all illusion.

The spirit is what we have witnessed manifested throughout the ages, but any age or historical record mean nothing because it is not in the present. Thus true presence is the unwavering spirit throughout the flow of the universe.

The Imperium is now, because it is the highest expression of the universal will for even when millions of decades pass our order will re-manifest into it's proper and divine peace will be known once more.

\hfill

\texttt{Logres on 2014-03-02 at 00:32 said:}

\url{http://www.millinerd.com/2011/01/horton-on-election-modern-purgatory.html}

``Not surprisingly, Hart sees the poverty of this perspective as a direct result of the loss of that old millinerd hobbyhorse, the analogia entis (the ``covenant of light''), which again, is shorthand for patristic/medieval metaphysical consensus, that transcendence -- the only Christian kind -- which transcends classical notions of transcendence. Christian ``being'' is almost always caricatured before it is deconstructed. But Xavier Zubiri explains the actual Christian view: ``Being is operation. And the more perfect something is, the deeper and more fertile is its operative activity. Being, said Dionysios, is ecstatic.'' But back to Hart:

What is absent from the Calvinist/Bañezian picture of divine causality is that ancient metaphysical vision that Przywara chose to call the `analogia entis'. In this ``analogical ontology', the infinite dependency of created being upon divine being is understood strictly in terms of the ever-greater difference between them; and, under the rule of this ontology, it is possible to affirm the real participation of the creature's freedom in God's free creative act without asserting any ontic continuity of kind between created and divine acts. When, however, the rule of analogy declines -- as it did at the threshold of modernity -- then invariably the words we attempt to apply both to creatures and to God (goodness, justice, mercy, love, freedom) dissolve into equivocity, and theology can recover its coherence only by choosing a single `attribute' to treat as univocal, in order that God and world might be united again. In the early modern period, the attribute most generally preferred was `power' or `sovereignty' -- or, more abstractly, 'cause.'

Calvin's theological errors, therefore, were the result of an atmospheric poverty. We trust the medievals not because they are old, but because of the air they breathed; and whether people realize it or not, people like C.S. Lewis because, as a medievalist, he breathed it too... ''

Notice that ``the Covenant of Light'' is the exact exoteric counterpart to the project of the Meditations, in Tomberg, and that the Great Chain of Being expresses this as a metaphysical Philosophy, the Queen of Sciences.

\hfill

\texttt{Michel on 2014-03-02 at 12:40 said:}

Lovely article Mr. Logres. I have some questions though, is the modern pagan thought and influence widespread? Also (regardless of whether or not it is widespread) , does it really pose what may be called a threat? Perhaps being a recluse I am unaware of certain things. Finally, at a certain point, doesn't the line between ``esoteric'' and ``exoteric'' become non-existent? Isn't it just One Savior? One Faith? One Spirit? One Love? (I know enough at least to sometimes still feel the effects of the misuse of those words.)

From the Infinite Archetypes, everything has been fashioned, God in His Wisdom has determined their measurements viz, limiting them and manifesting/creating them according to their own Archetypal Nature, their initial states as Limited Realities ``occuring'' as Being; the thing that is the sufficient reason for this creation being Beauty or Beatitude, the result of which is Love. There must be Beauty for there to be Love. But there must be Light for something akin to perception of Beauty to occur. ``The Light'' as a Reality can be seen on the one hand as God, or on the other hand as the ``intermediary'' necessary for Creation itself to occur. This is an ``intermediary'' without conjunction or mixture. It is simply an intermediary (and the first point of view, i.e. viewing it as nothing other than God, for His presence is never bereft of Radiance, finds its validation in this, no conjunction, no mixture, just Light). Of course this intermediary has been revealed to us as many names in the form of avatars and prophets, but it has been chiefly revealed as Jesus Christ. Moving on, the Presence of Christ enabled the Archetypes to see the Beauty of God. Seeing this Beauty, they fell in Love with Him and wanted this Love to continue forever. How could this be done? 

Digressing, when we see something we deem beautiful, we fall in love with it (for instance, a man is almost immediately attracted to a beautiful woman and can love her without even knowing her name), and seek to unite ourselves with it. But upon uniting ourselves with it, being one with it, there is nothing left to deem beautiful, to ``desire'' , to Love. The replacement of this void, is the ``desire'' for the perpetuation of the experience we had when we first united ourselves to it. The form of Love has changed therefore, from loving the thing, to loving the perpetuation of the experience of the thing. Thus we ``desire'' this and it can reach a point where the strength of this Love is so strong that we forget all else and are immersed in the experience; for instance, during the sexual experience at the epitome of the orgasm.

The perpetuation of this state of Love, could only be if God recreated the Universe with each instant/moment. So that, when the Universe, or rather, when we are recreated (by virtue of ``the intermediary of The Light'') with each instant/moment this recreation conforming to our Archetypal Natures, which is why it is called manifestation or unveiling- ``unveiling of our Archetypal Natures via limitation into Being'' we see the Beauty of God and fall in Love with Him all over again, and hence we are ``Being''. All the activities of all creatures are nothing other than, to a lesser or greater degree, a modulation of this Activity.The creatures reflect this Love, this ``thing'' that has caused them to ``BE'' in a haphazard way; they think they love this or that, but in truth, it is only God that they seek, it is only He that they Love. The process that brought them here must be ``overturned'' for them to realize this, they must see the essence of the traces; it is in this physical world that they are fully ``expressed'' as a Possibility, the salt has ``encapsulated the beings'' so that this ``fixity'' mimics their Immutable Archetypal Nature in its Changelessness, its ``Infinite Fixity''. Just as they were unveiled by Light to be manifested, so they must be unveiled by Light to see who they really Love, or they must be veiled by darkness just as God covers Himself in ``thick dark clouds''. What is called the ``exoteric'' is just a veiling of this Love; but Christ has torn the veil separating us from it, we can once more, from our ``fixed state as salt'' which mimics what we are as Immutable Archetypes via His ``intermediary'' as ``The Light'' see the Beauty of God and fall in Love Him not in the haphazard way characterized by the life we knew before , so that by virtue of this Love we are ``Being'' and as such, we can experience this Love, manifested as the creation of US as the Universe with each moment; in short, so that by this Grace, we can be saved. This is the Grace that overturns the cosmogonic process by repeating it in an upside down manner (by ascension as compared to descent or ``falling'').

So with that (please forgive me for my long comments) I fail to see the exoteric or esoteric, only the One Love of God and different ``measures'' of perceiving it; this Love has never changed and it ``endures forever''.

\hfill

\texttt{scardanelli on 2014-03-02 at 13:34 said:}

Thank you for the link to the CS Lewis quote, Logres... is was a warm refreshing breeze for the soul. Another man enamored of all things medieval, Kenelm Henry Digby has a fascinating little book called Maxims of Christian Chivalry- a collection from his larger work the Broad Stone of Honor. I've come back to it off and on when in need of inspiration. From the foreword: ``Nay! Nay! Chivalry is not dead, even in our day of sordid thought and selfish aim. It only slumbers. It will awaken, careless of all gain, defiant of all danger, devoted, impetuous, enthusiastic as ever crusader, when a man recognizes his vocation to personal honour... ''

\hfill

\texttt{Logres on 2014-03-03 at 02:39 said:}

Thanks, Scardanelli, for the reference.

Michel, I agree, but notwithstanding Borella and others' insistence that the essence of Christianity is ``the esoteric made exoteric'', the religion is based around paradoxes to begin with, so that esoteric and exoteric (being terms we use to denote spiritual experience) remain with their meaning, despite the eventual vision in which they merge in the Light. Just as there are a lot (or at least a few notable) pagans who have ``Christian bones'', but remain formally pagan, there are accurate spiritual contours denoted by the use of the terms you mentioned (exoteric, esoteric). Specifically, it's the manner in which ``God'', the ``Light'', or ``Love'' is experienced -- is it primarily through rituals or dogmas or practices, or does it flow through images that surface in the soul? There may be quite a few pagans farther along in Christian esoteric practice than many nominal Christians will ever hope to be. It's a strange, paradoxical, and uncomfortable fact (along with many others), but it seems to match with Divine and human truth. I don't understand why this is, but it is a mystery, and in order to comprehend it, one has to enter into the Arcanum. Since those who do this tend to be of higher spiritual quality than the rest of us, we call them ``saints'' or ``prophets'' or some other term, and speak of them having an esoteric truth that others lack, or that others can only formulate through art, poetry, or ritual, and experience in fleeting fashion, but hopefully more fully after death. That's the only way of harmonizing seemingly conflicting truths, in the paradox. Or, even better, through meditation on the paradox. Ideally, however, you are correct, \& perhaps we can say that the perfect man is the one for whom the esoteric and exoteric aspects are seamless realities, either entered at will, and accurately correlated to each other.

\hfill

\texttt{Logres on 2014-03-03 at 16:48 said:}

Michel, better answer here:

\url{http://www.gornahoor.net/?p=7152}

''Exoteric philosophy comprised ancient Greek philosophy. Esoteric philosophy is identical with the Christian religion''

Notice that in this syllogism, exoteric is identified with the esoteric Logos or inner Fire of the universe (of the Stoics), and that esoteric is identified with outward ritual. So this is a chiasm, in which earth echoes heaven, and heaven echoes earth. But yet (and this is the part where I may be differentiating a little from you) both remain more fully themselves, a paradox.

\hfill

\texttt{Michel on 2014-03-04 at 14:54 said:}

Thanks for the answer: ``the essence of Christianity is the esoteric made exoteric''.

``So this is a chiasm, in which earth echoes heaven, and heaven echoes earth. But yet (and this is the part where I may be differentiating a little from you) both remain more fully themselves, a paradox.''

I hope I didn't convey the notion that as modes of being, there is no distinction between heaven and earth and the subtle world too . I meant that after accomplishing the Hero's Journey, the World in its entirety is seen in a completely new way by the being that would have accomplished it. As modes of being, they each have their own conditions, but the Hero has gone beyond these conditions, the Hero has realized that within himself which acted as the sufficient reason for the prevailing of these conditions, and has conquered them.

``There may be quite a few pagans farther along in Christian esoteric practice than many nominal Christians will ever hope to be. It's a strange, paradoxical, and uncomfortable fact (along with many others), but it seems to match with Divine and human truth.''

I don't know if Freemasonry can still be considered as something akin to a Christian Path, but even in the preliminary stages for example, when one is locked with a skull and bones in a dark part of the lodge, preceding one's initiation, one meditates deeply on the futility of his prior perceptions of life (the room's setting acting as a support for the initial destruction of his false and profane self), and is ready, for the death that is to befall him, to re-begin a re-orientation of this perception of life, having being reduced to his materia prima following the utterance of the fiat lux; as he begins the initiatic voyage, his perceptions of the world merge from within, realizing and actualizing according to his aptitudes, the possibilities within himself, learning to live again building Solomon's Temple ; but you say this yourself, the paradox is resolved the deeper one sinks into the Arcanum. The stumbling block for most is the persistence of some of these perceptions even after they have, to a lesser or greater extent, began the path; they never truly mastered silence/death and hence began the path with their own pre-conceived notions, albeit subtle and ``crafty'' (in their solemnity and apparent reality), as the goal, thus they may end up worse than they were before, having gone deeper into a certain perception that they should have earlier destroyed e.g. a notion of what ``heaven is'',a bias either for or against something, a ``belief'' etc.

In fact, in the beginning, there should be no attempt at all to try and fit into something or to escape anything, first, it should be an acceptance of what one is, whatever that may be, followed by a seamless cathartic collapse or a sudden realization of the futility of what people call ``me''; something akin to a gestalt. One advantage of being under someone's tutelage, is that they can help with this deconstructive process; to use an Islamic concept to better explain it; there are certain Sheikhs or Muqadims (someone who is like a stand-in for a Sheikh) who exert something called the Himma. This is akin to the fiat lux, in that it is a ``reproduction'' of the Divine Command that created Order out of Chaos; the Tutor has the ability to directly assert his `will' towards the disciple so that the disciple is freed from certain difficulty. This can help the disciple in the beginning to completely destroy himself (his profane self) and achieve silence/initiatic death as well as later on, when the disciple is under a heavy burden, or is stuck at a certain point in the initiatic voyage. That means that by his Himma or you can say ``will-power'' (not forceful, just harmonious and in tandem with nature), a Tutor can help a disciple experience higher states of consciousness,literally- show him the essence of things even without much effort from the disciple. So the point is that perhaps the lingering of certain pre-conceived perceptions is what makes some Christians not able to make as much progress as pagans. (I hope by pagans you don't mean ``new-age'', I take it you mean philosophers like Plotinus and Aristotle).

\end{sffamily}
\end{footnotesize}

